% Curs 12 - Pagina 5 (3-I-1995).
% Fragment pentru % Curs 12 - Pagina 5 (3-I-1995).
% Fragment pentru % Curs 12 - Pagina 5 (3-I-1995).
% Fragment pentru % Curs 12 - Pagina 5 (3-I-1995).
% Fragment pentru \input{courses/pages/curs12-pagina5-ro}.
% Transcriere in ordinea manuscrisului; fara preambul si fara salt fortat de pagina.
% NOTE EDITORIALE: Formulele/propozitiile cu semnul intrebarii reflecta
% intrebarile din manuscris si nu sunt completate cu demonstratii inventate.

\subsection*{Curs 12 - Pagina 5}

\noindent\underline{Aplica\c{t}ie 1}\quad
``Monoid liber comutativ generat de o mul\c{t}ime''.

\[
R=\bigl\{(xy,yx)\mid x,y\in A^*\bigr\},\qquad R^{-1}=R.
\]

\medskip
\noindent\underline{Apl. 2}\quad Grup liber generat de o mul\c{t}ime.

\bigskip
\noindent\textbf{1.} Dac\u{a} $u,v\in A^*$, spunem c\u{a} $v$ se ob\c{t}ine
 din $u$ printr-o \underline{$R$-substitu\c{t}ie} dac\u{a}
\[
\exists\,(w,w')\in R\cup R^{-1}\cup\{(1,1)\},
\qquad \exists\,p,q\in A^*
\]
\noindent a.\^{\i}. 
\[
u=pwq,\qquad v=pw'q.
\]

\noindent (Adic\u{a} ele admit ca subcuvinte elementele unei perechi;
 au aceea\c{s}i parte ini\c{t}ial\u{a} \c{s}i final\u{a}.)

\bigskip
\noindent $S_R\subseteq A^*\times A^*$, rela\c{t}ie binar\u{a},
\[
u\,S_R\,v\quad\Longleftrightarrow\quad
v\text{ se ob\c{t}ine din }u\text{ printr-o }R\text{-substitu\c{t}ie}.
\]

\noindent $S_R$ e reflexiv\u{a}, simetric\u{a}, compatibil\u{a} la
 st\^{a}nga \c{s}i dreapta cu ``$\cdot$'' din $A^*$.
\[
u\,S_R\,v\ \overset{?}{\Longrightarrow}\
\begin{cases}
zu\,S_R\,zv, & (\forall)\,z\in A^*,\\
uz\,S_R\,vz, & (\forall)\,z\in A^*.
\end{cases}
\]
\noindent $zu=zpwq$ \c{s}i $zv=zpw'q$.

\medskip
\noindent $u\,S_R\,u\;?$ \quad $(\forall)\,u\in A^*$.
\[
u=u\cdot1\cdot1,\qquad u=u\cdot1\cdot1
\quad\bigl(w=w'=1\bigr).
\]

\noindent $u\,S_R\,v\ \overset{?}{\Longrightarrow}\ v\,S_R\,u$.

\bigskip
\noindent\underline{Obs.}\quad $R\subseteq S_R$.
\[
(w,w')\in R,\qquad p=q=1
\quad\Longrightarrow\quad (w,w')\in S_R.
\]

\bigskip
\noindent $R^*\subseteq A^*\times A^*$, unde, pentru $u,v\in A^*$,
\[
(u,v)\in R^*\ \Longleftrightarrow\
\begin{gathered}
\exists\,u_1,\ldots,u_n\in A^*,\quad n\text{ finit},\\
u=u_1,\quad v=u_n,\quad
u_i\,S_R\,u_{i+1}\quad (\forall)\,i=1,\ldots,n-1.
\end{gathered}
\]
% In manuscris, termenul de deasupra lui n este taiat si inlocuit cu „finit”.

\bigskip
\noindent\underline{Ex.}\quad $A=\{x,y,z\}$,
\[
R=\bigl\{(xy,yx),(yz,zy),(zx,xz)\bigr\}.
\]
\[
u=xxzyyx,\qquad v=yzxxyx.
\]
% VERIFICARE MANUSCRIS: Ultimele doua cuvinte sunt scrise compact.
% Citirile „xxzyyx” si „yzxxyx” sunt provizorii.
% Pagina se opreste la aceste exemple, fara o derivare suplimentara.
.
% Transcriere in ordinea manuscrisului; fara preambul si fara salt fortat de pagina.
% NOTE EDITORIALE: Formulele/propozitiile cu semnul intrebarii reflecta
% intrebarile din manuscris si nu sunt completate cu demonstratii inventate.

\subsection*{Curs 12 - Pagina 5}

\noindent\underline{Aplica\c{t}ie 1}\quad
``Monoid liber comutativ generat de o mul\c{t}ime''.

\[
R=\bigl\{(xy,yx)\mid x,y\in A^*\bigr\},\qquad R^{-1}=R.
\]

\medskip
\noindent\underline{Apl. 2}\quad Grup liber generat de o mul\c{t}ime.

\bigskip
\noindent\textbf{1.} Dac\u{a} $u,v\in A^*$, spunem c\u{a} $v$ se ob\c{t}ine
 din $u$ printr-o \underline{$R$-substitu\c{t}ie} dac\u{a}
\[
\exists\,(w,w')\in R\cup R^{-1}\cup\{(1,1)\},
\qquad \exists\,p,q\in A^*
\]
\noindent a.\^{\i}. 
\[
u=pwq,\qquad v=pw'q.
\]

\noindent (Adic\u{a} ele admit ca subcuvinte elementele unei perechi;
 au aceea\c{s}i parte ini\c{t}ial\u{a} \c{s}i final\u{a}.)

\bigskip
\noindent $S_R\subseteq A^*\times A^*$, rela\c{t}ie binar\u{a},
\[
u\,S_R\,v\quad\Longleftrightarrow\quad
v\text{ se ob\c{t}ine din }u\text{ printr-o }R\text{-substitu\c{t}ie}.
\]

\noindent $S_R$ e reflexiv\u{a}, simetric\u{a}, compatibil\u{a} la
 st\^{a}nga \c{s}i dreapta cu ``$\cdot$'' din $A^*$.
\[
u\,S_R\,v\ \overset{?}{\Longrightarrow}\
\begin{cases}
zu\,S_R\,zv, & (\forall)\,z\in A^*,\\
uz\,S_R\,vz, & (\forall)\,z\in A^*.
\end{cases}
\]
\noindent $zu=zpwq$ \c{s}i $zv=zpw'q$.

\medskip
\noindent $u\,S_R\,u\;?$ \quad $(\forall)\,u\in A^*$.
\[
u=u\cdot1\cdot1,\qquad u=u\cdot1\cdot1
\quad\bigl(w=w'=1\bigr).
\]

\noindent $u\,S_R\,v\ \overset{?}{\Longrightarrow}\ v\,S_R\,u$.

\bigskip
\noindent\underline{Obs.}\quad $R\subseteq S_R$.
\[
(w,w')\in R,\qquad p=q=1
\quad\Longrightarrow\quad (w,w')\in S_R.
\]

\bigskip
\noindent $R^*\subseteq A^*\times A^*$, unde, pentru $u,v\in A^*$,
\[
(u,v)\in R^*\ \Longleftrightarrow\
\begin{gathered}
\exists\,u_1,\ldots,u_n\in A^*,\quad n\text{ finit},\\
u=u_1,\quad v=u_n,\quad
u_i\,S_R\,u_{i+1}\quad (\forall)\,i=1,\ldots,n-1.
\end{gathered}
\]
% In manuscris, termenul de deasupra lui n este taiat si inlocuit cu „finit”.

\bigskip
\noindent\underline{Ex.}\quad $A=\{x,y,z\}$,
\[
R=\bigl\{(xy,yx),(yz,zy),(zx,xz)\bigr\}.
\]
\[
u=xxzyyx,\qquad v=yzxxyx.
\]
% VERIFICARE MANUSCRIS: Ultimele doua cuvinte sunt scrise compact.
% Citirile „xxzyyx” si „yzxxyx” sunt provizorii.
% Pagina se opreste la aceste exemple, fara o derivare suplimentara.
.
% Transcriere in ordinea manuscrisului; fara preambul si fara salt fortat de pagina.
% NOTE EDITORIALE: Formulele/propozitiile cu semnul intrebarii reflecta
% intrebarile din manuscris si nu sunt completate cu demonstratii inventate.

\subsection*{Curs 12 - Pagina 5}

\noindent\underline{Aplica\c{t}ie 1}\quad
``Monoid liber comutativ generat de o mul\c{t}ime''.

\[
R=\bigl\{(xy,yx)\mid x,y\in A^*\bigr\},\qquad R^{-1}=R.
\]

\medskip
\noindent\underline{Apl. 2}\quad Grup liber generat de o mul\c{t}ime.

\bigskip
\noindent\textbf{1.} Dac\u{a} $u,v\in A^*$, spunem c\u{a} $v$ se ob\c{t}ine
 din $u$ printr-o \underline{$R$-substitu\c{t}ie} dac\u{a}
\[
\exists\,(w,w')\in R\cup R^{-1}\cup\{(1,1)\},
\qquad \exists\,p,q\in A^*
\]
\noindent a.\^{\i}. 
\[
u=pwq,\qquad v=pw'q.
\]

\noindent (Adic\u{a} ele admit ca subcuvinte elementele unei perechi;
 au aceea\c{s}i parte ini\c{t}ial\u{a} \c{s}i final\u{a}.)

\bigskip
\noindent $S_R\subseteq A^*\times A^*$, rela\c{t}ie binar\u{a},
\[
u\,S_R\,v\quad\Longleftrightarrow\quad
v\text{ se ob\c{t}ine din }u\text{ printr-o }R\text{-substitu\c{t}ie}.
\]

\noindent $S_R$ e reflexiv\u{a}, simetric\u{a}, compatibil\u{a} la
 st\^{a}nga \c{s}i dreapta cu ``$\cdot$'' din $A^*$.
\[
u\,S_R\,v\ \overset{?}{\Longrightarrow}\
\begin{cases}
zu\,S_R\,zv, & (\forall)\,z\in A^*,\\
uz\,S_R\,vz, & (\forall)\,z\in A^*.
\end{cases}
\]
\noindent $zu=zpwq$ \c{s}i $zv=zpw'q$.

\medskip
\noindent $u\,S_R\,u\;?$ \quad $(\forall)\,u\in A^*$.
\[
u=u\cdot1\cdot1,\qquad u=u\cdot1\cdot1
\quad\bigl(w=w'=1\bigr).
\]

\noindent $u\,S_R\,v\ \overset{?}{\Longrightarrow}\ v\,S_R\,u$.

\bigskip
\noindent\underline{Obs.}\quad $R\subseteq S_R$.
\[
(w,w')\in R,\qquad p=q=1
\quad\Longrightarrow\quad (w,w')\in S_R.
\]

\bigskip
\noindent $R^*\subseteq A^*\times A^*$, unde, pentru $u,v\in A^*$,
\[
(u,v)\in R^*\ \Longleftrightarrow\
\begin{gathered}
\exists\,u_1,\ldots,u_n\in A^*,\quad n\text{ finit},\\
u=u_1,\quad v=u_n,\quad
u_i\,S_R\,u_{i+1}\quad (\forall)\,i=1,\ldots,n-1.
\end{gathered}
\]
% In manuscris, termenul de deasupra lui n este taiat si inlocuit cu „finit”.

\bigskip
\noindent\underline{Ex.}\quad $A=\{x,y,z\}$,
\[
R=\bigl\{(xy,yx),(yz,zy),(zx,xz)\bigr\}.
\]
\[
u=xxzyyx,\qquad v=yzxxyx.
\]
% VERIFICARE MANUSCRIS: Ultimele doua cuvinte sunt scrise compact.
% Citirile „xxzyyx” si „yzxxyx” sunt provizorii.
% Pagina se opreste la aceste exemple, fara o derivare suplimentara.
.
% Transcriere in ordinea manuscrisului; fara preambul si fara salt fortat de pagina.
% NOTE EDITORIALE: Formulele/propozitiile cu semnul intrebarii reflecta
% intrebarile din manuscris si nu sunt completate cu demonstratii inventate.

\subsection*{Curs 12 - Pagina 5}

\noindent\underline{Aplica\c{t}ie 1}\quad
``Monoid liber comutativ generat de o mul\c{t}ime''.

\[
R=\bigl\{(xy,yx)\mid x,y\in A^*\bigr\},\qquad R^{-1}=R.
\]

\medskip
\noindent\underline{Apl. 2}\quad Grup liber generat de o mul\c{t}ime.

\bigskip
\noindent\textbf{1.} Dac\u{a} $u,v\in A^*$, spunem c\u{a} $v$ se ob\c{t}ine
 din $u$ printr-o \underline{$R$-substitu\c{t}ie} dac\u{a}
\[
\exists\,(w,w')\in R\cup R^{-1}\cup\{(1,1)\},
\qquad \exists\,p,q\in A^*
\]
\noindent a.\^{\i}. 
\[
u=pwq,\qquad v=pw'q.
\]

\noindent (Adic\u{a} ele admit ca subcuvinte elementele unei perechi;
 au aceea\c{s}i parte ini\c{t}ial\u{a} \c{s}i final\u{a}.)

\bigskip
\noindent $S_R\subseteq A^*\times A^*$, rela\c{t}ie binar\u{a},
\[
u\,S_R\,v\quad\Longleftrightarrow\quad
v\text{ se ob\c{t}ine din }u\text{ printr-o }R\text{-substitu\c{t}ie}.
\]

\noindent $S_R$ e reflexiv\u{a}, simetric\u{a}, compatibil\u{a} la
 st\^{a}nga \c{s}i dreapta cu ``$\cdot$'' din $A^*$.
\[
u\,S_R\,v\ \overset{?}{\Longrightarrow}\
\begin{cases}
zu\,S_R\,zv, & (\forall)\,z\in A^*,\\
uz\,S_R\,vz, & (\forall)\,z\in A^*.
\end{cases}
\]
\noindent $zu=zpwq$ \c{s}i $zv=zpw'q$.

\medskip
\noindent $u\,S_R\,u\;?$ \quad $(\forall)\,u\in A^*$.
\[
u=u\cdot1\cdot1,\qquad u=u\cdot1\cdot1
\quad\bigl(w=w'=1\bigr).
\]

\noindent $u\,S_R\,v\ \overset{?}{\Longrightarrow}\ v\,S_R\,u$.

\bigskip
\noindent\underline{Obs.}\quad $R\subseteq S_R$.
\[
(w,w')\in R,\qquad p=q=1
\quad\Longrightarrow\quad (w,w')\in S_R.
\]

\bigskip
\noindent $R^*\subseteq A^*\times A^*$, unde, pentru $u,v\in A^*$,
\[
(u,v)\in R^*\ \Longleftrightarrow\
\begin{gathered}
\exists\,u_1,\ldots,u_n\in A^*,\quad n\text{ finit},\\
u=u_1,\quad v=u_n,\quad
u_i\,S_R\,u_{i+1}\quad (\forall)\,i=1,\ldots,n-1.
\end{gathered}
\]
% In manuscris, termenul de deasupra lui n este taiat si inlocuit cu „finit”.

\bigskip
\noindent\underline{Ex.}\quad $A=\{x,y,z\}$,
\[
R=\bigl\{(xy,yx),(yz,zy),(zx,xz)\bigr\}.
\]
\[
u=xxzyyx,\qquad v=yzxxyx.
\]
% VERIFICARE MANUSCRIS: Ultimele doua cuvinte sunt scrise compact.
% Citirile „xxzyyx” si „yzxxyx” sunt provizorii.
% Pagina se opreste la aceste exemple, fara o derivare suplimentara.
